\emph{Översikt}
Så länge vi själva bestämde vad programmen räknade på räknade de alltid
rätt.  När användaren matar in värdena går det inte längre: en del av
det hon kan skriva går inte att räkna med.  Den här modulen handlar om
vad programmet ska göra då.  Vi utgår från ett program som frågar efter
ett födelseår och ser vad som händer när svaret inte är ett årtal ---
körningen \emph{avbryts}, och Python skriver en \emph{spårutskrift}
(\foreignlanguage{english}{traceback}).  Vi
lär oss läsa den, och sedan ta hand om felet själva med
\mintinline{python}{try} och \mintinline{python}{except}: en gren per
sorts fel, de specifika först och det allmänna sist eller inte alls,
och med ett meddelande skrivet för användaren i stället för för
programmeraren.  Sist ser vi hur våra egna funktioner kan \emph{kasta}
särfall vidare till den som anropat dem.

\emph{Lärandemål}
Efter denna modul ska studenten kunna:
% Lo04 (styrstrukturer), Lo08 (interaktiva program)
\begin{restatable}{lo}{ExcLOException}\label{ExcLOException}%
  Förklara vad ett särfall (\foreignlanguage{english}{exception}) är,
  och redogöra för vad som händer när ett särfall inträffar.
\end{restatable}
% Lo08 (interaktiva program), Lo11 (granska program)
\begin{restatable}{lo}{ExcLOTraceback}\label{ExcLOTraceback}%
  Läsa en spårutskrift (\foreignlanguage{english}{traceback}) och ur
  den utläsa var felet uppstod, vilken sorts
  fel det var och vad som var fel.
\end{restatable}
% Lo04 (styrstrukturer), Lo08 (interaktiva program)
\begin{restatable}{lo}{ExcLOCatch}\label{ExcLOCatch}%
  Fånga särfall med \mintinline{python}{try} och
  \mintinline{python}{except} --- de specifika sorterna före den
  allmänna --- och redogöra för vad som händer i programmet.
\end{restatable}
% Lo05 (användbara utskrifter), Lo08 (interaktiva program)
\begin{restatable}{lo}{ExcLOMessage}\label{ExcLOMessage}%
  Utforma ett felmeddelande som användaren kan agera på, i stället för
  att låta programmet krascha eller vidarebefordra Pythons egen text.
\end{restatable}
% Lo03 (dela upp program), Lo04 (styrstrukturer)
\begin{restatable}{lo}{ExcLORaise}\label{ExcLORaise}%
  Läsa och skriva \mintinline{python}{raise} i en egen funktion, och
  förklara hur ett särfall sprider sig uppåt till den som anropat
  funktionen.
\end{restatable}

\emph{Förkunskaper}
Modulen förutsätter variabler, datatyperna \mintinline{python}{int},
\mintinline{python}{float} och \mintinline{python}{str}, utskrifter och
egna funktioner med parametrar och returvärden från tidigare, samt
inmatning med \mintinline{python}{input} och typomvandling från
föreläsningen \emph{Inmatning och datatyper}, och
\mintinline{python}{while} från föreläsningen \emph{Villkor och
styrstrukturer}.

\emph{Läsning}
Pythons egen dokumentation beskriver mekanismen i kapitlet \emph{Errors
and Exceptions}.\footnote{\label{fn:pydoc}%
  \url{https://docs.python.org/3/tutorial/errors.html} (hämtad
  2026-09-03).}
